\documentclass{umk-matematika}

\begin{document}

\maketitlepage
    {Практическая работа №13}
    {Теорема о трёх перпендикулярах}
    {}

\section{Фундамент (1--4 балла)}

\textbf{Задача 1 (1 балл).} Сформулируйте прямую теорему о трёх перпендикулярах.

\textbf{Задача 2 (2 балла).} Сформулируйте обратную теорему о трёх перпендикулярах.

\textbf{Задача 3 (3 балла).} Из точки A к плоскости α проведены перпендикуляр AB и наклонная AC. Что является проекцией наклонной AC на плоскость α?

\textbf{Задача 4 (4 балла).} В треугольнике ABC угол C прямой. Отрезок CD — перпендикуляр к плоскости ABC. Назовите проекцию наклонной AD на плоскость ABC.

\section{Стены (5--7 баллов)}

\textbf{Задача 5 (5 баллов).} Из точки M к плоскости α проведён перпендикуляр MO = 8 см и наклонная MA = 17 см. Проекция наклонной OA = 15 см. Перпендикулярна ли прямая a, проведённая в плоскости через точку A к OA, наклонной MA?

\textbf{Задача 6 (6 баллов).} В кубе $ABCDA_1B_1C_1D_1$ докажите, что прямая BD перпендикулярна прямой AC₁.

\textbf{Задача 7 (7 баллов).} Из точки A к плоскости проведён перпендикуляр AB = 9 см и наклонная AC = 15 см. В плоскости через точку C проведена прямая, перпендикулярная BC. Найдите расстояние от A до этой прямой.

\section{Кровля (8--10 баллов)}

\textbf{Задача 8 (8 баллов).} Из вершины A квадрата ABCD проведён перпендикуляр AK к плоскости квадрата. Сторона квадрата = 6 см, AK = 8 см. Найдите расстояние от точки K до прямой BD.

\textbf{Задача 9 (9 баллов).} Из точки M к плоскости α проведены перпендикуляр MO и две наклонные MA и MB. Известно, что MA = 10 см, MB = 17 см, а их проекции относятся как 6:15. Найдите расстояние от M до плоскости.

\textbf{Задача 10 (10 баллов).} Строительная мачта высотой 12 м укреплена тросом, закреплённым на земле на расстоянии 5 м от основания мачты. Перпендикулярна ли линия крепления троса к земле самой мачте, если трос натянут под прямым углом к линии на земле, соединяющей основание мачты с точкой крепления?

\newpage
\section*{Ответы}

\begin{enumerate}
  \item прямая теорема.
  \item обратная теорема.
  \item отрезок BC.
  \item отрезок AC.
  \item да, перпендикулярна.
  \item BD ⟂ AC₁.
  \item 15 см.
  \item $\sqrt{82}$ см.
  \item $MO = 8$ см.
  \item да, перпендикулярна (мачта ⟂ земле, линия ⊂ земле).
\end{enumerate}

\end{document}